\section{Practical Recommendations}
\label{sec:recs}

The finding is only useful if it changes what practitioners do. These
recommendations follow directly from the failure mode above and do not depend on
resolving which judge is correct --- they are precautions warranted by the
observed hazard.

\begin{enumerate}
\item \textbf{Do not let a model be the primary judge of its own family.} In our
data the largest real signal (per an independent judge) sat in exactly the
outputs the self-family judge floored to zero variance. If you must score a
model's outputs, use a judge from a different family --- ideally a stronger one
--- for at least a validation subsample. Treating a same-family judge's score as
independent ground truth is the specific mistake this paper documents.

\item \textbf{Grep for zero-variance cells.} A judge that returns the same score
for every item in a condition (or for an entire subject model) cannot detect a
within-condition effect, and this is silent: the pipeline succeeds, the numbers
look clean, and the null is an artifact of the instrument's floor, not the
behavior. Flagging any (judge, subject, dimension) cell with zero or near-zero
variance is a cheap, code-only check that would have caught this before any
conclusion was drawn.

\item \textbf{Cross-check against a stronger judge on a subsample, and report
$\kappa$.} A single frontier-judge pass over a stratified subsample (here, 180
items cost \$3.36) is enough to expose a conclusion-flipping disagreement.
Report between-judge agreement per dimension: aggregate agreement computed over
easy floor/ceiling dimensions can look reassuring while agreement on the
discriminating dimensions is at chance.

\item \textbf{Do not assume a multi-judge ensemble is independent.} Averaging
three judges that share a training lineage (or simply a small model's compressed
dynamic range) can reproduce a blind spot rather than cancel it. Check that
panel members actually disagree on the dimensions that matter before trusting
their consensus.

\item \textbf{Add a cheap non-model audit as a tripwire, not an oracle.} A
deterministic keyword/structure check is coarse, but when it agrees with a
suspiciously floored judge and \emph{disagrees} with a stronger judge, that gap
is a useful signal that the disagreement is about subtle, contextual behavior ---
exactly where small judges are least reliable.

\item \textbf{Report the judge as part of the claim.} ``Model X was not more
manipulative under pressure'' is not a fact about Model X; it is a fact about
Model X \emph{as scored by judge Y}. State the judge, its relationship to the
subject, and the correction procedure alongside every LLM-judged effect.
\end{enumerate}

These are not aspirational. The committed \texttt{pressure\_evals.judge\_diagnostics}
module implements the zero-variance / dynamic-range check (recommendations 1--2)
and regenerates the judge$\times$subject dynamic-range table
(Table~\ref{tab:dynamic_range}) and the self-family comparison directly from the
frozen scores; the between-judge agreement statistics (recommendation 3) and the
deterministic audit (recommendation 5) are likewise first-class, reproducible
steps. Running the dynamic-range check on our own data is what surfaced that the
primary judge floored every subject model --- the signal that a single
same-family number would have hidden.
